% Source: https://artofproblemsolving.com/community/c3h1197839_3d_geometry
% Stable identity: aops-topic-1197839-q01
Note that $s^3=4\cdot2\sqrt2=8\sqrt2$.

\emph{Regular tetrahedron:} volume $\dfrac{s^3}{6\sqrt2}=\dfrac{8\sqrt2}{6\sqrt2}=\dfrac{4}{3}$; two of them give $\dfrac{8}{3}$.

\emph{Corner tetrahedron:} the three right isosceles faces meet at a corner $O$ with three mutually perpendicular edges of length $a=\dfrac{s}{\sqrt2}$ (so the base has side $a\sqrt2=s$).

\begin{asy}
usepackage("amsmath");
size(3cm);
defaultpen(fontsize(10pt));
triple A=(0,0,0),B=(1,0,0),C=(0,1,0),D=(0,0,1);
pair p(triple P){ return (P.x+0.45*P.y, P.z+0.3*P.y); }
draw(p(A)--p(B)--p(D)--cycle); draw(p(B)--p(C)--p(D)); draw(p(A)--p(C),dashed);
label("$O$",p(A),SW); label("$a$",(p(A)+p(B))/2,S); label("$a$",(p(A)+p(D))/2,W); label("$a$",(p(A)+p(C))/2,NW);
label("$a\sqrt2$",(p(B)+p(D))/2,NE);
label("\mbox{}",(min(currentpicture,user=true)+max(currentpicture,user=true))/2,black); // keeps the text after the figure black when LaTeX typesets the labels inline
\end{asy}

Its volume is $\dfrac{a^3}{6}=\dfrac{s^3}{12\sqrt2}=\dfrac{2}{3}$; two give $\dfrac{4}{3}$.

\emph{Square pyramid with all edges $s$:} the height is $\dfrac{s}{\sqrt2}$, so the volume is $\dfrac{1}{3}s^2\cdot\dfrac{s}{\sqrt2}=\dfrac{s^3}{3\sqrt2}=\dfrac{8}{3}$.

\emph{Last tetrahedron:} two equilateral faces $ABC$ and $ABD$ share the edge $AB$, and the right isosceles faces $ACD$, $BCD$ force $CD=s\sqrt2$. This is exactly half of the square pyramid above: cutting that pyramid (apex $P$, base $ABCD$) by the plane through $P$, $A$, $C$ gives the tetrahedron $PABC$, whose faces $PAB$, $PBC$ are equilateral and whose faces $ABC$, $PAC$ are right isosceles with hypotenuse $AC=s\sqrt2$. So its volume is $\dfrac{1}{2}\cdot\dfrac{8}{3}=\dfrac{4}{3}$.

Total: $\dfrac{8}{3}+\dfrac{4}{3}+\dfrac{8}{3}+\dfrac{4}{3}=\boxed{8}$.
